\documentclass[11pt]{article}

\usepackage[margin=1in]{geometry}
\usepackage{amsmath, amssymb}
\usepackage{booktabs}
\usepackage{enumitem}
\usepackage{graphicx}
\usepackage{hyperref}
\usepackage{longtable}
\usepackage{tabularx}
\usepackage{comment}
\usepackage{docmute}
\graphicspath{{../../artifacts/showcase/plots/}{../../artifacts/bayesian_walkthroughs/plots/}{../../artifacts/transition_matrix_accumulator_v1/}{../../artifacts/corpus_autodevelopment_v1/plots/}{../../artifacts/generic_corpus_exploration/}}

\title{Kinematic Classifier Methodology Compendium}
\author{kinematic-classifier-sandbox}
\date{\today}

\begin{document}

\maketitle

\begin{abstract}
This compendium joins the current workflow, classifier-ladder, and corpus-search notes into one structured reference. It is intended to be the broad technical front door for the repository: it explains how a study is declared, how a corpus is generated and audited, how evidence-provider families construct posteriors, and how promotion decisions are justified. The current 1D witness suite remains a proof of methodology layers rather than the final scientific endpoint. Its surrounding method vocabulary is closest to Bayesian filtering, reinforcement-learning-style search, and quality-diversity exploration rather than to a single monolithic classifier claim \cite{barshalom,bishop,sutton_barto,mouret_clune}.
\end{abstract}

\tableofcontents
\clearpage

\section{How To Read This Compendium}
This document is organized around three complementary questions:
\begin{enumerate}[leftmargin=1.5em]
  \item How does a study move through the repository end to end?
  \item How do classifier and filter rungs construct evidence and earn promotion?
  \item How does the corpus layer generate, score, and defend a selected dataset?
\end{enumerate}

The three parts below intentionally preserve the companion-note split:
\begin{itemize}[leftmargin=1.5em]
  \item the workflow part explains the study object and end-to-end flow,
  \item the ladder part explains evidence contracts and escalation logic,
  \item the corpus part explains candidate generation, adequacy, and selection.
\end{itemize}

\section{Scope And Claims}
This compendium makes a methodological claim, not a final domain claim. The strongest current statement is:
\begin{quote}
\textit{The repository already demonstrates a reusable study workflow for kinematic classification in which corpus governance, evidence construction, posterior updates, and promotion decisions are made under explicit contracts.}
\end{quote}

That claim is intentionally narrower than any claim such as ``the final classifier architecture is complete'' or ``the current 1D witness suite is the target scientific benchmark.'' The current witness suite is still used as a methodology proof:
\begin{itemize}[leftmargin=1.5em]
  \item they show how a study is declared,
  \item they show how evidence-provider families are compared under a shared posterior contract,
  \item they show how corpus adequacy can be scored and audited,
  \item they do not yet exhaust the space of higher-dimensional kinematic backends or final production corpora.
\end{itemize}

The generalization story therefore runs in one direction:
\begin{equation}
\text{1D witness success}
\Longrightarrow
\text{workflow and contract plausibility},
\end{equation}
but not yet
\begin{equation}
\text{1D witness success}
\Longrightarrow
\text{all future 3D studies are solved}.
\end{equation}

\section{Notation Conventions}
The same symbols recur across all three parts:
\begin{center}
\begin{tabularx}{0.92\textwidth}{@{}lX@{}}
\toprule
Symbol & Meaning \\
\midrule
\(\mathcal{D}\) or \(\mathcal{D}^{\star}\) & corpus, dataset slice, or selected corpus candidate \\
\(f\) or \(\mathcal{F}\) & feature set or active feature family \\
\(\mathcal{C}\) & class set or class pair under study \\
\(m\) & classifier, filter, or evidence-provider family \\
\(\pi\) & prior regime \\
\(b\) & backend, sensing, or dynamics family \\
\(z_t\) & raw observation at time \(t\) \\
\(\phi_t\) & derived feature summary at time \(t\) \\
\(\ell_t(c)\) & class log-evidence contribution at time \(t\) \\
\(p_t(c)\) & posterior class probability after time \(t\) \\
\(d\) & study-level decision in \(\{\text{promote}, \text{revise}, \text{reject}, \text{defer}\}\) \\
\bottomrule
\end{tabularx}
\end{center}

Two naming choices are used consistently in this compendium:
\begin{itemize}[leftmargin=1.5em]
  \item \textbf{Corpus Explorer} refers to the generic corpus-generation, selection, and adequacy machinery.
  \item \textbf{Study Candidate Evaluator} refers to the generic study-level evaluator that consumes a selected corpus, feature set, class set, method family, and prior regime to produce a promotion decision.
\end{itemize}

\section{Reader Guidance}
The parts can be read independently, but they are intended to support one larger argument:
\begin{enumerate}[leftmargin=1.5em]
  \item Part I defines the end-to-end study object and explains where each repository surface fits.
  \item Part II explains how evidence-provider families share one posterior contract and how a stronger rung earns promotion.
  \item Part III explains how a corpus candidate becomes credible enough to support the Part II comparisons at all.
\end{enumerate}

\part{Workflow}
\section*{Part I Bridge}
\addcontentsline{toc}{section}{Part I Bridge}
Part I is the front door. It defines the study object, the corpus layer, the evidence layer, and the decision layer before the method-specific details appear. Read it first if the main question is ``what overall workflow is this repository trying to prove out?''

\begingroup
\let\maketitle\relax
\excludecomment{abstract}
\documentclass[11pt]{article}

\usepackage[margin=1in]{geometry}
\usepackage{amsmath, amssymb}
\usepackage{booktabs}
\usepackage{enumitem}
\usepackage{graphicx}
\usepackage{hyperref}
\usepackage{longtable}
\usepackage{tabularx}
\graphicspath{{../../artifacts/showcase/plots/}{../../artifacts/bayesian_walkthroughs/plots/}}

\title{Kinematic Classifier Workflow and Math}
\author{kinematic-classifier-sandbox}
\date{\today}

\begin{document}

\maketitle

\begin{abstract}
This document is a practical walkthrough of the repository's kinematic-classification workflow. Its purpose is narrower than the full methodology paper: it explains what flows through the system, how corpus, feature, evidence, and posterior layers connect, and which equations govern the main study loop. The current 1D witness suite is used as a readable proof of workflow layers, not as the final scientific endpoint. The workflow language combines corpus governance, Bayesian posterior updating, and evidence-provider comparisons using standard statistical references as the mathematical baseline rather than claiming a new end-to-end theory \cite{bishop2006,kalman1960,brier1950}.
\end{abstract}

\section{Purpose}
The repository can be read at two levels:
\begin{enumerate}[leftmargin=1.5em]
  \item as a generic study workflow for kinematic classification, and
  \item as a current 1D witness suite that proves parts of that workflow.
\end{enumerate}

This document focuses on the workflow itself. The main question is:
\begin{quote}
\textit{How does a declared kinematic classification study move from corpus design to a promotion decision under explicit mathematical contracts?}
\end{quote}

\section{End-To-End Flow}
The shortest accurate summary of the repository is
\begin{align}
\text{study declaration}
&\rightarrow
\text{corpus generation or selection}
\rightarrow
\text{class-validity and coverage audits}
\notag \\
&\rightarrow
\text{feature extraction or model-based evidence}
\rightarrow
\text{posterior updates}
\notag \\
&\rightarrow
\text{evaluation, stress tests, and promotion decision}.
\end{align}

The central study object is
\begin{equation}
s = (\mathcal{D}, f, \mathcal{C}, m, \pi, b),
\end{equation}
where
\begin{align}
\mathcal{D} &:= \text{corpus or selected dataset slice}, \\
f &:= \text{feature set}, \\
\mathcal{C} &:= \text{class set or class pair}, \\
m &:= \text{classifier or filter family}, \\
\pi &:= \text{prior regime}, \\
b &:= \text{optional backend or dynamics family}.
\end{align}

The workflow then has two generic engines:
\begin{align}
\text{Corpus Explorer} &: (o, b) \mapsto \mathcal{D}^{\star}, \\
\text{Study Candidate Evaluator} &: (\mathcal{D}^{\star}, f, \mathcal{C}, m, \pi, b) \mapsto d,
\end{align}
with decision
\begin{equation}
d \in \{\text{promote}, \text{revise}, \text{reject}, \text{defer}\}.
\end{equation}

The current strongest supported corpus-evaluation workflow is the default generated/common-study path. The selected-corpus rerun path is also strong because the chosen generated packet is pushed back through adequacy, leakage, and regression checks. An adapter/contract layer exists for provided corpus sources, but it is not yet a full arbitrary-corpus adequacy pipeline, so the repo should not be described as providing full generic corpus evaluation for any provided corpus.

Figure~\ref{fig:study_object_map} is the compact front-door diagram for that object.

\begin{figure}[htbp]
\centering
\setlength{\fboxsep}{8pt}
\fbox{
\begin{minipage}{0.92\textwidth}
\centering
\begin{tabularx}{\textwidth}{@{}XcXcX@{}}
\textbf{Declared Study} &
\(\rightarrow\) &
\textbf{Corpus Engine} &
\(\rightarrow\) &
\textbf{Decision Engine} \\
\(s = (\mathcal{D}, f, \mathcal{C}, m, \pi, b)\) &
&
\((o, b) \mapsto \mathcal{D}^{\star}\) &
&
\((\mathcal{D}^{\star}, f, \mathcal{C}, m, \pi, b) \mapsto d\) \\
\end{tabularx}

\vspace{0.7em}

\begin{tabularx}{\textwidth}{@{}>{\raggedright\arraybackslash}X>{\raggedright\arraybackslash}X>{\raggedright\arraybackslash}X@{}}
\textbf{\(\mathcal{D}\)} corpus slice or selected dataset &
\textbf{\(f\)} feature set and evidence summary family &
\textbf{\(\mathcal{C}\)} class set or class pair \\
\textbf{\(m\)} classifier or filter family &
\textbf{\(\pi\)} prior regime &
\textbf{\(b\)} backend, sensor, or dynamics family \\
\end{tabularx}
\end{minipage}
}
\caption{The repository's front-door study object. The left box declares what is being claimed, the middle engine selects a corpus that can test the claim honestly, and the right engine decides whether the declared study is promotable.}
\label{fig:study_object_map}
\end{figure}

\section{Workflow Stages}
\subsection{Stage 1: Declare the Study}
The workflow starts by declaring what is being tested rather than jumping directly to a model run. That declaration fixes:
\begin{itemize}[leftmargin=1.5em]
  \item the classes to separate,
  \item the kinematic feature family to use,
  \item the evidence-provider family to compare,
  \item the prior regime,
  \item and the corpus requirements needed to make the claim honest.
\end{itemize}

This separation matters because poor results can have multiple causes:
\begin{equation}
\text{failure} \neq \text{classifier failure only}.
\end{equation}
A bad result can indicate inadequate corpus coverage, weak features, ambiguous class definitions, prior fragility, or a missing algorithmic rung.

\subsection{Stage 2: Generate or Select a Corpus}
For the current 1D witness suite, a trajectory is represented as
\begin{equation}
\tau_i = \{(t_{ik}, z_{ik})\}_{k=1}^{n_i},
\end{equation}
with latent kinematic state
\begin{equation}
x_{ik} =
\begin{bmatrix}
p_{ik} \\
v_{ik} \\
a_{ik}
\end{bmatrix},
\qquad
z_{ik} = h(x_{ik}) + \epsilon_{ik}.
\end{equation}

In the generic story, the corpus layer defines a distribution over scenarios rather than one fixed benchmark file:
\begin{equation}
\tau \sim p(\tau \mid \theta, g, \eta),
\end{equation}
where \(\theta\) are scenario parameters, \(g\) is a class or motion template, and \(\eta\) captures noise, corruption, or stress settings.

The corpus objective is not merely ``make more data.'' It is to produce a dataset that exercises the intended boundaries:
\begin{align}
Q_{\text{corpus}}
=\;&
w_{\text{balance}} Q_{\text{balance}}
+ w_{\text{pair}} Q_{\text{pair}}
+ w_{\text{feature}} Q_{\text{feature}}
\notag \\
&+ w_{\text{difficulty}} Q_{\text{difficulty}}
- w_{\text{leak}} P_{\text{leak}}
- w_{\text{invalid}} P_{\text{invalid}}.
\end{align}
The exact score terms vary by analysis surface, but the structure is stable: reward coverage and difficulty diversity, penalize leakage and invalid examples.

\subsection{Stage 3: Audit Coverage Before Classifier Claims}
Before comparing classifiers, the workflow checks whether the corpus itself is methodologically usable. The main audit questions are:
\begin{itemize}[leftmargin=1.5em]
  \item Do declared classes remain valid after generation and corruption?
  \item Are feature sets actually excited by the generated trajectories?
  \item Are hard class pairs represented near their intended boundaries?
  \item Is the corpus balanced enough to support downstream interpretation?
  \item Are there obvious covariate leaks that make the task artificially easy?
\end{itemize}

In compact form, the adequacy gate is
\begin{equation}
\text{corpus adequate}
\Longleftrightarrow
\bigwedge_{j=1}^{J} G_j(\mathcal{D}) = \text{pass},
\end{equation}
where each \(G_j\) is a gate such as class validity, balance, feature-set coverage, or leakage.

The canonical artifact path for this stage is a paired adequacy-and-coverage packet: \texttt{corpus\_adequacy\_audit\_v1} establishes the corpus gates themselves, while \texttt{coverage\_report\_v1} records classifier-support and feature-set readiness against the same audited corpus.

\subsection{Stage 4: Convert Observations Into Evidence}
The repository separates raw observations from derived features and from posterior-ready evidence:
\begin{align}
z_k &= \text{raw observation at time } t_k, \\
\phi_k &= \Phi(z_{1:k}, t_{1:k}), \\
y_k &= \text{evidence-provider input at step } k.
\end{align}

The evidence input \(y_k\) depends on the method family:
\begin{equation}
y_k =
\begin{cases}
z_k, & \text{pointwise methods}, \\
\phi_k, & \text{windowed feature methods}, \\
\nu_{k,c}, & \text{Kalman innovation evidence}, \\
(\phi_k, \nu_{k,c}), & \text{hybrid or future combined methods}.
\end{cases}
\end{equation}

The important design rule is that methods are compared through a shared evidence contract:
\begin{equation}
\mathcal{E}(z_{1:k}, \phi_{1:k}, \text{model state}) \rightarrow \{\ell_k(c)\}_{c \in \mathcal{C}},
\end{equation}
where \(\ell_k(c)\) is the class log-evidence contribution at time \(k\).

\subsection{Stage 5: Update the Posterior}
Once evidence is expressed in a compatible form, the posterior update is generic. Let
\begin{equation}
p_{k}(c) := p(c \mid y_{1:k}).
\end{equation}
Then the recursive Bayes update is
\begin{equation}
p_k(c) \propto p_{k-1}(c) \exp(\ell_k(c)),
\end{equation}
with normalized form
\begin{equation}
p_k(c) =
\frac{p_{k-1}(c) \exp(\ell_k(c))}
{\sum_{c' \in \mathcal{C}} p_{k-1}(c') \exp(\ell_k(c'))}.
\end{equation}

In log space, define
\begin{equation}
\alpha_k(c) = \log p_{k-1}(c) + \ell_k(c),
\end{equation}
then
\begin{equation}
p_k(c) = \frac{\exp(\alpha_k(c))}
{\sum_{c' \in \mathcal{C}} \exp(\alpha_k(c'))}.
\end{equation}

For a two-class problem with classes \(c_1\) and \(c_2\), the log-odds recursion is
\begin{equation}
\log \frac{p_k(c_1)}{p_k(c_2)}
=
\log \frac{p_{k-1}(c_1)}{p_{k-1}(c_2)}
+
\ell_k(c_1) - \ell_k(c_2).
\end{equation}
This equation is central to the repo because it exposes prior sensitivity and flip thresholds directly.

\subsection{Stage 6: Add Structure When Static Evidence Is Not Enough}
If pointwise or windowed evidence is insufficient, the workflow adds structure one rung at a time rather than leaping immediately to the most complex filter.

\paragraph{Transition-aware accumulation.}
For switching-aware inference, class posteriors are propagated through a transition matrix \(T\):
\begin{equation}
\tilde{p}_k(c') = \sum_{c \in \mathcal{C}} T_{cc'} \, p_{k-1}(c),
\end{equation}
followed by the evidence update
\begin{equation}
p_k(c') \propto \tilde{p}_k(c') \exp(\ell_k(c')).
\end{equation}

\paragraph{Model-based innovation evidence.}
For Kalman-bank style methods, each class carries a dynamics model and emits an innovation \(\nu_{k,c}\) with covariance \(S_{k,c}\). The class evidence increment is
\begin{equation}
\ell_k(c)
=
-\frac{1}{2}
\left(
\nu_{k,c}^{\top} S_{k,c}^{-1} \nu_{k,c}
+ \log \det S_{k,c}
+ d \log(2\pi)
\right),
\end{equation}
where \(d\) is the observation dimension.

The ladder logic is therefore:
\begin{align}
\text{pointwise}
\rightarrow
\text{windowed}
\rightarrow
\text{sequential Bayes}
\rightarrow
\text{Kalman bank}
\notag \\
\rightarrow
\text{transition-aware filters}
\rightarrow
\text{advanced filters}.
\end{align}
Each rung must fix a demonstrated failure mode of the previous one.

\subsection{Stage 7: Evaluate the Study}
After posterior histories are produced, the evaluator does more than inspect top-line accuracy. It measures:
\begin{itemize}[leftmargin=1.5em]
  \item confusion and pairwise failure localization,
  \item prior sensitivity,
  \item separability of the feature space,
  \item calibration and posterior sharpness,
  \item robustness under stress, corruption, and boundary cases.
\end{itemize}

The evaluation layer asks two different questions:
\begin{align}
\text{Is the feature/class geometry promising?} \\
\text{Does the chosen evidence-provider family realize that promise?}
\end{align}

That is why the workflow includes both static geometry diagnostics and dynamic posterior diagnostics.

\subsection{Stage 8: Promote, Revise, Reject, or Defer}
The final output is a decision, not just a score:
\begin{equation}
d = \Delta(
\text{corpus audits},
\text{separability},
\text{posterior quality},
\text{stress results}
).
\end{equation}

The point of the decision layer is to keep the repository honest. A study should not be promoted because one artifact looked good in isolation. It should be promoted only when the corpus, feature, evidence, and posterior stories are mutually consistent.

\section{Worked Numeric Examples}
This section uses current repository artifacts so the equations above can be checked against real runs rather than only against symbolic notation.

\subsection{Sequential Bayes Example}
The simplest hand-checkable update in the current artifacts is the binary demo recorded in the Bayes step table artifact. It starts with a uniform prior
\begin{equation}
p_0(A) = p_0(B) = 0.5
\end{equation}
and one-step log evidence values
\begin{align}
\ell_1(A) &= -0.2231435513, \\
\ell_1(B) &= -1.6094379124.
\end{align}

The log Bayes factor is therefore
\begin{equation}
\ell_1(A) - \ell_1(B) = 1.3862943611.
\end{equation}
With equal priors, the posterior odds equal the Bayes factor, so
\begin{equation}
\log \frac{p_1(A)}{p_1(B)} = 1.3862943611
\end{equation}
and the normalized posterior is
\begin{equation}
\bigl(p_1(A), p_1(B)\bigr) = (0.8, 0.2).
\end{equation}

The same artifact also contains a real common-study trajectory where the final log posterior odds for
\texttt{constant\_velocity} versus \texttt{constant\_acceleration} reach
\begin{equation}
\log \frac{p(\texttt{constant\_velocity} \mid y_{1:4})}
{p(\texttt{constant\_acceleration} \mid y_{1:4})}
= 0.7005868961,
\end{equation}
giving final posteriors
\begin{equation}
\bigl(0.6683178817,\; 0.3316821183\bigr).
\end{equation}

\subsection{Kalman Innovation Example}
For the Kalman-bank rung, the scalar innovation likelihood for class \(c\) is
\begin{equation}
\ell_k(c)
=
-\frac{1}{2}
\left(
\frac{\nu_{k,c}^2}{S_{k,c}}
+ \log S_{k,c}
+ \log(2\pi)
\right).
\end{equation}

Using the row for \texttt{constant\_velocity\_regular\_0} at step \(4\) from
\texttt{innovation\_history.csv}, the
\texttt{constant\_velocity} model has
\begin{align}
\nu_{k,\mathrm{cv}} &= -0.1727524922, \\
S_{k,\mathrm{cv}} &= 0.1350588971.
\end{align}
Substitution gives
\begin{equation}
\ell_k(\mathrm{cv})
=
-\frac{1}{2}
\left(
\frac{(-0.1727524922)^2}{0.1350588971}
+ \log(0.1350588971)
+ \log(2\pi)
\right)
= -0.028399.
\end{equation}

On the same step, the \texttt{constant\_acceleration} model has
\begin{align}
\nu_{k,\mathrm{ca}} &= 0.4632798272, \\
S_{k,\mathrm{ca}} &= 0.4137106782,
\end{align}
which yields
\begin{equation}
\ell_k(\mathrm{ca}) = -0.737038.
\end{equation}
Those exact values match the corresponding posterior-history row, where the
\texttt{constant\_velocity} model dominates that step.

\subsection{Transition-Matrix Example}
The transition-aware numeric walkthrough
shows the propagation step explicitly. At the first switched step for the target
mode \texttt{braking},
\begin{equation}
\bar{p}_t(\texttt{braking})
= \sum_s p_{t-1}(s)\,T_{s,\texttt{braking}}
= 0.312.
\end{equation}

The contributing terms are
\begin{align}
0.340 \cdot 0.800 &= 0.272, \\
0.660 \cdot 0.060 &= 0.040,
\end{align}
with the remaining source-mode contributions rounding to zero, so the propagated prior mass becomes
\begin{equation}
0.272 + 0.040 = 0.312.
\end{equation}

That walkthrough also records the empirical effect:
\begin{itemize}[leftmargin=1.5em]
  \item static post-switch accuracy: \(0.000\),
  \item transition-matrix post-switch accuracy: \(0.333\).
\end{itemize}
So the key mechanism is not a new emission model, but moving prior mass into a plausible switched mode before the emission term is applied.

\subsection{Corpus Adequacy Example}
The corpus-audit scorecard in
\texttt{corpus\_adequacy\_report.md}
currently reports
\begin{align}
Q_{\text{corpus}} &= 0.643, \\
L &= 0.613, \\
T &= 0.500, \\
G &= 0.023.
\end{align}

The same scorecard shows the positive coverage terms
\begin{align}
B_{\text{class}} &= 1.000, \\
B_{\text{tier}} &= 1.000, \\
B_{\text{covariates}} &= 0.704, \\
E_{\text{feature}} &= 0.557, \\
C_{\text{pair}} &= 0.730, \\
V &= 1.000.
\end{align}

The gate still fails even though feature coverage passes, because one hard pair is too easy:
\begin{equation}
\texttt{constant\_acceleration} \text{ vs. } \texttt{maneuver}
\end{equation}
has observed
\begin{align}
\text{pairwise AUC} &= 1.000, \\
\text{overlap} &= 0.000,
\end{align}
despite being declared \texttt{hard}. This is exactly the sort of evidence the audit is meant to surface: a ``good'' classifier result on that pair would be untrustworthy because the corpus boundary is not yet credible.

\subsection{Advanced-Filter Decision Example}
The escalation gate for stronger switching or nonlinear filters is also evidence-based. The walkthrough in
\texttt{advanced\_filter\_decision\_numeric\_walkthrough.md}
reports
\begin{align}
\Delta_{\mathrm{post-switch}}^{\mathrm{TM-static}} &= 0.102, \\
\Delta_{\mathrm{overall}}^{\mathrm{TM-static}} &= 0.046, \\
\Delta_{\mathrm{post-switch}}^{\mathrm{TM-kalman}} &= 0.151, \\
\Delta_{\mathrm{overall}}^{\mathrm{TM-kalman}} &= 0.123.
\end{align}

The current implemented decisions are therefore
\begin{align}
\text{IMM justified now} &= \text{False}, \\
\text{particle filter justified now} &= \text{False}.
\end{align}
The important methodological point is that the repo does not promote a stronger evidence-provider family merely because it exists. It promotes that family only after a simpler rung fails for a reason the next rung can specifically explain.

\section{Visual Anchors}
The numeric examples above are easier to internalize when paired with the existing artifact plots. Figure~\ref{fig:bayes_transition_visuals} shows the two most direct posterior-oriented examples, and Figure~\ref{fig:governance_visuals} shows the governance layer that decides whether stronger claims are justified.

\begin{figure}[htbp]
\centering
\includegraphics[width=0.48\textwidth]{posterior_timeline.png}
\hfill
\includegraphics[width=0.48\textwidth]{transition_matrix_diagnostics.png}
\caption{Left: posterior evolution for the Bayes walkthrough family. Right: the transition-matrix witness showing how switching-aware propagation changes behavior after mode changes.}
\label{fig:bayes_transition_visuals}
\end{figure}

\begin{figure}[htbp]
\centering
\includegraphics[width=0.48\textwidth]{corpus_adequacy_scorecard.png}
\hfill
\includegraphics[width=0.48\textwidth]{advanced_filter_decision_matrix.png}
\caption{Left: corpus governance is evaluated before classifier claims are promoted. Right: escalation to more complex filters is decision-gated by measured witness evidence rather than by ambition alone.}
\label{fig:governance_visuals}
\end{figure}

\section{Implementation Map}
Table~\ref{tab:workflow_map} ties the abstract stages to current code surfaces and typical artifacts.

\small
\begin{longtable}{@{}p{0.16\textwidth} p{0.25\textwidth} p{0.27\textwidth} p{0.24\textwidth}@{}}
\caption{Current workflow stages, code surfaces, and representative artifacts.}
\label{tab:workflow_map} \\
\toprule
Stage & Main question & Code surfaces & Representative artifacts \\
\midrule
\endfirsthead
\toprule
Stage & Main question & Code surfaces & Representative artifacts \\
\midrule
\endhead
Study declaration & What classes, features, priors, and method family are being claimed? & \texttt{common\_experiment}; study-candidate helpers; manifests in \texttt{experiments/} & study manifests; feature sets; class-pair manifests \\
Corpus generation & What trajectories should be generated or selected? & \texttt{trajectory\_generator}; \texttt{corpus/exploration} & selected corpus manifests; candidate score tables \\
Validity and adequacy & Is the corpus methodologically usable? & \texttt{class\_validity}; \texttt{coverage\_report} & \texttt{corpus\_adequacy\_audit\_v1/}; balance and leakage reports \\
Feature geometry & Are the classes separated in the active feature space? & \texttt{feature\_analysis}; \texttt{pca\_analysis} & \texttt{feature\_analysis\_v1/}; \texttt{pca\_analysis\_v1/} \\
Evidence provider & How is class evidence constructed at each step? & sequential Bayes; Kalman bank; \texttt{transition\_matrix} & posterior histories; innovation tables; transition diagnostics \\
Promotion logic & Is the declared study promotable? & \texttt{validation\_ladder}; reporting surfaces & validation-ladder reports; promotion scorecards \\
\bottomrule
\end{longtable}
\normalsize

\section{What This Document Proves And What It Does Not}
This document proves the current workflow shape and the core mathematical contracts:
\begin{itemize}[leftmargin=1.5em]
  \item corpus, feature, evidence, posterior, and decision layers are explicitly separated,
  \item method families are compared through shared log-evidence and posterior contracts,
  \item the 1D witness suite is a readable proof of workflow layers,
  \item the repository already has places to audit adequacy before classifier claims.
\end{itemize}

This document does not prove:
\begin{itemize}[leftmargin=1.5em]
  \item that the current 1D corpus is final,
  \item that one method family is universally best,
  \item that 3D lift is complete,
  \item or that every conceptual contract is fully generalized across all backends today.
\end{itemize}

The honest current claim is:
\begin{quote}
\textit{The repository already supports a coherent kinematic-classification workflow with explicit mathematical interfaces, and the 1D witness suite is there to prove those interfaces before broader dimensional and backend lift.}
\end{quote}

\section{Artifact Guide}
Table~\ref{tab:artifact_guide} lists the artifact names referenced in this document. The full paths live under the corresponding artifact directories in the repository.

\small
\begin{table}[htbp]
\centering
\caption{Artifact names referenced directly in this workflow note.}
\label{tab:artifact_guide}
\begin{tabularx}{\textwidth}{@{}lX@{}}
\toprule
Artifact name & Purpose in this document \\
\midrule
\texttt{bayesian\_step\_tables.csv} & Hand-checkable and common-study Bayes posterior updates \\
\texttt{innovation\_history.csv} & Per-step Kalman innovations and innovation variances \\
\texttt{posterior\_history.csv} & Per-step posterior values used to confirm the likelihood calculations \\
\texttt{transition\_matrix\_numeric\_walkthrough.md} & Worked switching-aware propagation example \\
\texttt{corpus\_adequacy\_report.md} & Corpus scorecard and failure recommendations \\
\texttt{advanced\_filter\_decision\_numeric\_walkthrough.md} & Evidence-based gate for IMM and particle-filter escalation \\
\bottomrule
\end{tabularx}
\end{table}
\normalsize

\section{References}
\begin{thebibliography}{9}
\bibitem{bishop2006}
C. M. Bishop.
\newblock \emph{Pattern Recognition and Machine Learning}.
\newblock Springer, 2006.

\bibitem{kalman1960}
R. E. Kalman.
\newblock A new approach to linear filtering and prediction problems.
\newblock \emph{Transactions of the ASME}, 82D:35--45, 1960.

\bibitem{brier1950}
G. W. Brier.
\newblock Verification of forecasts expressed in terms of probability.
\newblock \emph{Monthly Weather Review}, 78(1):1--3, 1950.
\end{thebibliography}

\end{document}

\endgroup

\part{Classifier Ladder}
\section*{Part II Bridge}
\addcontentsline{toc}{section}{Part II Bridge}
Part II narrows from repository workflow to evidence math. It explains how pointwise, windowed, sequential, Kalman-style, and switching-aware methods are compared through one shared log-evidence and posterior-update grammar. Read it second if the main question is ``when is a stronger classifier rung actually justified?''

\begingroup
\let\maketitle\relax
\excludecomment{abstract}
\documentclass[11pt]{article}

\usepackage[margin=1in]{geometry}
\usepackage{amsmath, amssymb}
\usepackage{booktabs}
\usepackage{enumitem}
\usepackage{graphicx}
\usepackage{hyperref}
\usepackage{longtable}
\usepackage{tabularx}
\graphicspath{{../../artifacts/showcase/plots/}{../../artifacts/transition_matrix_accumulator_v1/}}

\title{Classifier Ladder Math}
\author{kinematic-classifier-sandbox}
\date{\today}

\begin{document}

\maketitle

\begin{abstract}
This document is the math companion to the workflow note. Its job is narrower than the full survey document: it states the shared evidence contract, shows how each ladder rung constructs class evidence, and explains the promotion rule for escalating from simpler to stronger evidence-provider families. The ladder vocabulary is grounded in standard Bayesian filtering, linear-Gaussian state estimation, and switching/nonlinear evidence constructions rather than a repo-specific probability calculus \cite{kalman1960,barshalom2001,gordon1993,doucet2001}.
\end{abstract}

\section{Purpose}
The classifier ladder answers one recurring question:
\begin{quote}
\textit{Given the observations available up to time \(t\), what is the next class-evidence increment and why is this rung justified?}
\end{quote}

The repository does not treat the ladder as a complexity ranking. It treats it as an evidence-design ladder:
\begin{equation}
\text{new rung} \Longrightarrow \text{new evidence capability} \Longrightarrow \text{new proof burden}.
\end{equation}

\section{Shared Contract}
Across the current ladder, the output contract is:
\begin{align}
y_t &:= \text{method-specific evidence input}, \\
\ell_t(c) &:= \text{class log-evidence at time } t, \\
\alpha_t(c) &= \log p_{t-1}(c) + \ell_t(c), \\
p_t(c) &= \frac{\exp(\alpha_t(c))}{\sum_j \exp(\alpha_t(j))}.
\end{align}

Each rung differs only in how it constructs \(\ell_t(c)\). The downstream posterior, calibration, confusion, and promotion surfaces are meant to stay shared.

\begin{table}[htbp]
\centering
\caption{The shared ladder contract.}
\label{tab:ladder_contract}
\begin{tabularx}{\textwidth}{@{}lXX@{}}
\toprule
Rung & Evidence input \(y_t\) & Evidence increment \(\ell_t(c)\) \\
\midrule
Pointwise & current observation or local measurement summary & \(\log p(y_t \mid c)\) \\
Windowed & fixed-window or prefix feature vector \(\phi_t\) & \(\log p(\phi_t \mid c)\) \\
Sequential Bayes & recursively emitted per-step evidence stream & cumulative additive log evidence \\
Kalman bank & innovation residual \(\nu_{t,c}\) with covariance \(S_{t,c}\) & \(\log \mathcal{N}(\nu_{t,c}; 0, S_{t,c})\) \\
Transition matrix & emitted evidence plus switching prior & \(\log \bar{p}_t(c) + \log E_t(c)\) \\
Advanced filters & mode-conditioned or particle-weighted evidence & method-specific switching or nonlinear evidence increment \\
\bottomrule
\end{tabularx}
\end{table}

\section{Ladder View}
The current ladder can be summarized as
\begin{align}
\text{pointwise}
\rightarrow
\text{windowed}
\rightarrow
\text{sequential Bayes}
\notag \\
\rightarrow
\text{Kalman bank}
\rightarrow
\text{transition matrix}
\rightarrow
\text{advanced filters}.
\end{align}

\begin{table}[htbp]
\centering
\caption{Current reader-facing ladder interpretation.}
\label{tab:ladder_story}
\begin{tabularx}{\textwidth}{@{}c l X X l@{}}
\toprule
Rung & Algorithm & Adds & Failure addressed & Witness \\
\midrule
0 & Pointwise & local likelihood baseline & no audited instantaneous baseline & pointwise overlap \\
1 & Windowed & short-history feature evidence & local noise and outlier fragility & windowed outlier/extrema \\
2 & Sequential Bayes & recursive memory & pointwise history blindness & sequential history \\
3 & Kalman bank & dynamics-conditioned innovation evidence & endpoint ambiguity under irregular timing & Kalman endpoint match \\
4 & Transition matrix & explicit switching pressure & static-class assumption & transition switching \\
5+ & Advanced filters & switching, nonlinear, or mixed-mode escalation & failures that simpler rungs cannot explain away & gated witnesses \\
\bottomrule
\end{tabularx}
\end{table}

\section{Pointwise Rung}
The pointwise rung asks whether a single observation already carries enough signal. With class \(c\), observation \(z_t\), and class parameters \((\mu_c, \sigma_c)\), the Gaussian pointwise evidence is
\begin{equation}
\ell_t(c)
=
-\frac{1}{2}
\left[
\log(2\pi \sigma_c^2)
+
\frac{(z_t - \mu_c)^2}{\sigma_c^2}
\right].
\end{equation}

The rung is justified only as a lower bound. If it works well, the study may be simple. If it fails, that does not yet justify a state-space model; it only proves that purely local evidence is insufficient.

\section{Windowed Rung}
The windowed rung changes the evidence atom from a raw observation to a feature vector
\begin{equation}
\phi_t = \Phi(z_{1:t}, t_{1:t}),
\end{equation}
with examples such as extrema, range, slope, sign changes, and residual summaries. Under the current diagonal-Gaussian approximation,
\begin{equation}
\ell_t(c)
=
\sum_{f \in \mathcal{F}}
\log \mathcal{N}\!\left(\phi_{t,f}; \mu_{c,f}, \sigma_{c,f}^2\right).
\end{equation}

The proof burden at this rung is not only accuracy. It must also show that the chosen features are the right answer to the previous failure mode, for example outlier sensitivity or short-horizon ambiguity.

\section{Sequential Bayes Rung}
The sequential Bayes rung keeps the same notion of per-step evidence but accumulates it recursively:
\begin{equation}
\log \frac{p_t(c_1)}{p_t(c_2)}
=
\log \frac{p_{t-1}(c_1)}{p_{t-1}(c_2)}
+
\ell_t(c_1) - \ell_t(c_2).
\end{equation}

This rung is justified when the pointwise or windowed method has the right local evidence direction but lacks enough temporal memory to turn weak repeated evidence into a stable decision.

\section{Kalman-Bank Rung}
The Kalman-bank rung replaces feature-engineered evidence with model-based innovation evidence. For class \(c\),
\begin{align}
\hat{x}_{t|t-1}^{(c)} &= F_t^{(c)} \hat{x}_{t-1|t-1}^{(c)}, \\
\nu_{t,c} &= z_t - H_t^{(c)} \hat{x}_{t|t-1}^{(c)}, \\
S_{t,c} &= H_t^{(c)} P_{t|t-1}^{(c)} H_t^{(c)\top} + R_t^{(c)}.
\end{align}
The resulting evidence increment is
\begin{equation}
\ell_t(c)
=
-\frac{1}{2}
\left(
\nu_{t,c}^{\top} S_{t,c}^{-1} \nu_{t,c}
+ \log \det S_{t,c}
+ d \log(2\pi)
\right).
\end{equation}

This rung is justified only when the previous failure is really about dynamics, timing, or state evolution. It is not justified merely because Kalman filters are more sophisticated.

\section{Transition-Matrix Rung}
The transition-matrix rung keeps the emission model but changes the prior propagation:
\begin{equation}
\bar{p}_t(c') = \sum_c T_{cc'} p_{t-1}(c).
\end{equation}
The posterior update becomes
\begin{equation}
p_t(c')
\propto
\bar{p}_t(c') \exp(\ell_t(c')).
\end{equation}

This rung is the bridge from static-class accumulation to explicit switching pressure. It is justified when the observed failure is mode switching, not merely noisy evidence.

\begin{figure}[htbp]
\centering
\includegraphics[width=0.8\textwidth]{transition_matrix_diagnostics.png}
\caption{Transition-matrix witness diagnostics. This rung should only be promoted when switching-aware propagation fixes a failure the static accumulator cannot.}
\label{fig:transition_matrix_ladder}
\end{figure}

\section{Promotion Rule}
The ladder promotion rule is methodological rather than aesthetic:
\begin{equation}
\text{promote rung } r+1
\Longleftrightarrow
\text{rung } r \text{ fails on a witness for a reason } r+1 \text{ is designed to fix.}
\end{equation}

That can be unfolded into four checks:
\begin{enumerate}[leftmargin=1.5em]
  \item The corpus and class geometry are credible enough that the failure is not a data-governance artifact.
  \item The previous rung fails on a controlled witness case.
  \item The next rung adds exactly the missing evidence capability.
  \item The improvement is visible in shared posterior and evaluation artifacts.
\end{enumerate}

This is the key discipline that prevents the ladder from turning into ``try a more complex model and hope.''

\section{Advanced-Filter Gate}
The advanced-filter families belong after the transition rung, not before it. Their proof burden is stricter:
\begin{itemize}[leftmargin=1.5em]
  \item IMM requires a switching witness case that simpler transition logic no longer explains adequately.
  \item PF requires a nonlinear or non-Gaussian failure that sensing improvements and robust Gaussian methods do not remove.
  \item RBPF requires a mixed discrete/continuous structure that neither plain PF nor plain Kalman-style switching handles cleanly.
\end{itemize}

The repository already encodes this as an evidence-based gate rather than an aspirational roadmap.

\begin{figure}[htbp]
\centering
\includegraphics[width=0.8\textwidth]{advanced_filter_decision_matrix.png}
\caption{Advanced-filter escalation should be gated by witness-suite evidence, not by general preference for complex algorithms.}
\label{fig:advanced_filter_gate}
\end{figure}

\section{What This Note Is For}
This document is intended to sit behind the workflow note:
\begin{itemize}[leftmargin=1.5em]
  \item the workflow note explains how a study moves through the repository,
  \item this ladder note explains how each method family constructs evidence and earns promotion.
\end{itemize}

The honest current claim is therefore limited but strong:
\begin{quote}
\textit{The repository already supports a shared posterior contract across multiple evidence-provider families, and it has a methodology for deciding when the next rung is actually justified.}
\end{quote}

\section{References}
\begin{thebibliography}{9}
\bibitem{kalman1960}
R. E. Kalman.
\newblock A new approach to linear filtering and prediction problems.
\newblock \emph{Transactions of the ASME}, 82D:35--45, 1960.

\bibitem{barshalom2001}
Y. Bar-Shalom, X. R. Li, and T. Kirubarajan.
\newblock \emph{Estimation with Applications to Tracking and Navigation}.
\newblock Wiley, 2001.

\bibitem{gordon1993}
N. J. Gordon, D. J. Salmond, and A. F. M. Smith.
\newblock Novel approach to nonlinear/non-Gaussian Bayesian state estimation.
\newblock \emph{IEE Proceedings F}, 140(2):107--113, 1993.

\bibitem{doucet2001}
A. Doucet, N. de Freitas, and N. Gordon, editors.
\newblock \emph{Sequential Monte Carlo Methods in Practice}.
\newblock Springer, 2001.
\end{thebibliography}

\end{document}

\endgroup

\part{Corpus Search}
\section*{Part III Bridge}
\addcontentsline{toc}{section}{Part III Bridge}
Part III moves back upstream and asks whether the dataset itself deserves trust. It explains candidate generation, adequacy gates, scalar scoring, Pareto tradeoffs, leakage penalties, and coverage reporting. Read it third if the main question is ``how do we know the classifier comparison is being run on an honest corpus?''

\begingroup
\let\maketitle\relax
\excludecomment{abstract}
\documentclass[11pt]{article}

\usepackage[margin=1in]{geometry}
\usepackage{amsmath, amssymb}
\usepackage{booktabs}
\usepackage{enumitem}
\usepackage{graphicx}
\usepackage{hyperref}
\usepackage{longtable}
\usepackage{tabularx}
\graphicspath{{../../artifacts/showcase/plots/}{../../artifacts/corpus_autodevelopment_v1/plots/}{../../artifacts/generic_corpus_exploration/}}

\title{Corpus Search Math}
\author{kinematic-classifier-sandbox}
\date{\today}

\begin{document}

\maketitle

\begin{abstract}
This document is the corpus-side math companion to the workflow note. It focuses on how corpus candidates are proposed, scored, audited, compared, and selected before classifier claims are trusted. Its optimization and archive vocabulary is intentionally borrowed from multiobjective search and quality-diversity work, with adaptive stress used as the motivating exploration pattern rather than as a claim of a new corpus optimizer \cite{deb2002,mouret2015,corso2022}.
\end{abstract}

\section{Purpose}
The corpus side of the repository answers a different question than the classifier ladder:
\begin{quote}
\textit{How do we generate or select a corpus that exercises the intended boundaries without making the task trivial, biased, or poorly governed?}
\end{quote}

The output object is not just ``more data.'' It is a selected corpus
\begin{equation}
\mathcal{D}^{\star}
\end{equation}
that is credible enough to support a declared study candidate.

\section{Corpus Explorer Contract}
The corpus engine is
\begin{equation}
(o, b, q, G_b) \mapsto \mathcal{D}^{\star},
\end{equation}
where
\begin{align}
o &:= \text{corpus objective}, \\
b &:= \text{backend or execution family}, \\
q &:= \text{candidate proposal distribution}, \\
G_b &:= \text{backend-specific trajectory generator}.
\end{align}

The selected corpus then feeds the study object
\begin{equation}
s = (\mathcal{D}^{\star}, f, \mathcal{C}, m, \pi, b).
\end{equation}
This is the reason corpus governance matters: a study candidate can fail before any classifier is run if \(\mathcal{D}^{\star}\) is not honest.

\section{Candidate Generation}
The generator side is a distribution over trajectory parameters rather than one fixed dataset:
\begin{align}
\theta_k &\sim q(\theta \mid o, b), \\
\tau_i &\sim G_b(\theta_k, \xi_i), \\
\mathcal{D}_k &= \{\tau_i\}_{i=1}^{N_k}.
\end{align}

The parameter block can be read as
\begin{equation}
\theta =
(\theta_{\text{class}},
\theta_{\text{tier}},
\theta_{\text{noise}},
\theta_{\text{sampling}},
\theta_{\text{env}}).
\end{equation}

That gives a compact story:
\begin{equation}
\text{objective}
\rightarrow
\text{candidate parameters}
\rightarrow
\text{generated corpus}
\rightarrow
\text{audit and selection}.
\end{equation}

\section{Adequacy Score}
The repo already evaluates corpora on balance, pair coverage, feature excitation, leakage, triviality, and degeneracy. A useful compact score form is
\begin{equation}
S_k
=
w_b B_k
+ w_c C_k
+ w_f F_k
+ w_d D_k
- w_\ell L_k
- w_t T_k
- w_g G_k,
\end{equation}
where
\begin{align}
B_k &:= \text{balance score}, \\
C_k &:= \text{class-pair coverage score}, \\
F_k &:= \text{feature excitation score}, \\
D_k &:= \text{difficulty-diversity score}, \\
L_k &:= \text{covariate leakage penalty}, \\
T_k &:= \text{triviality penalty}, \\
G_k &:= \text{degeneracy penalty}.
\end{align}

\begin{table}[htbp]
\centering
\caption{Operational meaning of the main corpus score terms.}
\label{tab:corpus_terms}
\begin{tabularx}{\textwidth}{@{}lXc@{}}
\toprule
Term & Meaning & Desired direction \\
\midrule
\(B_k\) & similarity to target class and tier balance & high \\
\(C_k\) & declared hard-pair boundary coverage & high \\
\(F_k\) & excitation of active feature sets & high \\
\(D_k\) & match to target difficulty distribution & high \\
\(L_k\) & class-linked shortcut signal in covariates & low \\
\(T_k\) & penalty for corpora that make hard pairs too easy & low \\
\(G_k\) & duplicate, invalid, or degenerate structure penalty & low \\
\bottomrule
\end{tabularx}
\end{table}

\section{Worked Score Example}
The current autodevelopment walkthrough gives one concrete example:
\begin{align}
B_k &= 1.000, \\
C_k &= 0.800, \\
F_k &= 0.720, \\
D_k &= 0.800, \\
L_k &= 0.102, \\
T_k &= 0.037, \\
G_k &= 0.000.
\end{align}
So for the selected candidate,
\begin{equation}
S_k
=
1.000 + 0.800 + 0.720 + 0.800 - 0.102 - 0.037 - 0.000
= 3.181.
\end{equation}

The important methodological point is that this scalar score is not the only decision surface. It is a summary. The repo also preserves the underlying audit terms and the non-dominated alternatives.

\section{Adequacy Gates}
The score alone is not sufficient. The corpus must also pass or at least clearly report adequacy gates:
\begin{equation}
\text{corpus adequate}
\Longleftrightarrow
\bigwedge_{j=1}^{J} G_j(\mathcal{D}_k) = \text{pass}.
\end{equation}

Typical gates include:
\begin{itemize}[leftmargin=1.5em]
  \item class validity,
  \item class balance,
  \item feature-set coverage,
  \item class-pair boundary coverage,
  \item covariate leakage,
  \item degeneracy.
\end{itemize}

One current audit example is especially important: a pair can be declared \texttt{hard} and still fail the governance check if the generated corpus makes it artificially easy. In that case, a strong classifier result is not evidence of a strong study.

\section{Pareto View}
Scalar ranking is useful for selection, but it can hide tradeoffs. The repo therefore also keeps a Pareto view using an objective vector
\begin{equation}
\mathbf{o}_k
=
\bigl[
B_k,\,
C_k,\,
F_k,\,
D_k,\,
-L_k,\,
-T_k,\,
-G_k
\bigr].
\end{equation}

Candidate \(a\) dominates candidate \(b\) when
\begin{align}
\mathbf{o}_a &\succeq \mathbf{o}_b \text{ componentwise}, \\
\mathbf{o}_a &\neq \mathbf{o}_b.
\end{align}

This matters because two corpora may have similar scalar scores but very different failure shapes. One may be slightly leakier; another may under-cover the intended hard boundaries. The Pareto front keeps both visible.

\begin{figure}[htbp]
\centering
\includegraphics[height=0.34\textheight,keepaspectratio]{corpus_score_pareto.png}
\caption{Corpus candidates should be compared as tradeoff points, not only as one scalar ranking. The Pareto view keeps non-dominated alternatives visible for study governance.}
\label{fig:corpus_pareto}
\end{figure}

\section{Exploration Reward}
The exploration surfaces go beyond one-shot ranking. A generic exploration reward can be read as
\begin{equation}
R_k
=
\lambda_1 U_{\text{coverage}}
+ \lambda_2 U_{\text{difficulty}}
+ \lambda_3 U_{\text{novelty}}
- \lambda_4 P_{\text{leak}}
- \lambda_5 P_{\text{invalid}},
\end{equation}
where the utility terms reward filling useful coverage gaps and the penalties reject shortcuts or invalid trajectories.

The important repository claim is not that this is a learned RL policy already. The claim is that the exploration surface is formalized enough to compare candidate-generation strategies through a common reward grammar.

\section{Coverage Reporting}
The corpus side now has a coverage-oriented reporting layer as well. It summarizes the adequacy gates, feature-set coverage, feature-group behavior, and classifier-support implications into one inspection bundle.

That means corpus evaluation can be read in two passes:
\begin{enumerate}[leftmargin=1.5em]
  \item the adequacy score and gate status,
  \item the detailed failure recommendations that explain what to fix next.
\end{enumerate}

\begin{figure}[htbp]
\centering
\includegraphics[width=0.72\textwidth]{corpus_adequacy_scorecard.png}
\caption{Corpus scorecards summarize whether the current candidate is usable at all before classifier results are trusted.}
\label{fig:corpus_scorecard}
\end{figure}

\begin{figure}[htbp]
\centering
\includegraphics[width=0.72\textwidth]{class_pair_coverage_heatmap.png}
\caption{Class-pair coverage diagnostics show whether the intended hard boundaries are actually being exercised.}
\label{fig:corpus_governance}
\end{figure}

\section{Leakage And Selection Risk}
The most important anti-pattern on the corpus side is shortcut leakage. If covariates such as duration, sampling irregularity, or noise level reveal the class too directly, then the corpus has stopped testing the intended kinematic distinction.

That is why a covariate audit is part of the corpus math story:
\begin{equation}
\text{good classifier result on leaky corpus}
\not\Rightarrow
\text{good study}.
\end{equation}

\begin{figure}[htbp]
\centering
\includegraphics[width=0.7\textwidth]{covariate_leakage_audit.png}
\caption{Covariate leakage is treated as a first-class penalty because it can invalidate downstream classifier conclusions even when accuracy looks strong.}
\label{fig:covariate_leakage}
\end{figure}

\section{What This Note Is For}
This note is meant to pair with the workflow and ladder notes:
\begin{itemize}[leftmargin=1.5em]
  \item the workflow note explains the end-to-end study loop,
  \item the ladder note explains how evidence-provider families earn promotion,
  \item this corpus note explains how the selected corpus is generated, judged, and defended.
\end{itemize}

The honest current claim is:
\begin{quote}
\textit{The repository already has a mathematically legible corpus-governance layer that can propose, audit, compare, and select corpora before study-level promotion decisions are made.}
\end{quote}

\section{References}
\begin{thebibliography}{9}
\bibitem{deb2002}
K. Deb, A. Pratap, S. Agarwal, and T. Meyarivan.
\newblock A fast and elitist multiobjective genetic algorithm: NSGA-II.
\newblock \emph{IEEE Transactions on Evolutionary Computation}, 6(2):182--197, 2002.

\bibitem{mouret2015}
J.-B. Mouret and J. Clune.
\newblock Illuminating search spaces by mapping elites.
\newblock \emph{arXiv}:1504.04909, 2015.

\bibitem{corso2022}
J. Corso, R. Moss, and M. Koren.
\newblock Adaptive stress testing and failure search in simulation.
\newblock \emph{Annual Review of Control, Robotics, and Autonomous Systems}, 5:59--83, 2022.
\end{thebibliography}

\end{document}

\endgroup

\clearpage
\section{Generalization Path And Current Limits}
The repository is not intended to stop at a toy 1D classifier exercise. Its broader claim is that the \emph{contracts} should generalize even when the current witnesses remain simple. The main reusable contracts are:
\begin{itemize}[leftmargin=1.5em]
  \item the study object \(s = (\mathcal{D}, f, \mathcal{C}, m, \pi, b)\),
  \item the Corpus Explorer contract for generating, auditing, and selecting \(\mathcal{D}^{\star}\),
  \item the shared evidence-provider contract that emits \(\ell_t(c)\),
  \item the shared posterior update and study-level promotion decision.
\end{itemize}

Those contracts should survive several kinds of lift:
\begin{enumerate}[leftmargin=1.5em]
  \item from 1D trajectories to higher-dimensional kinematic state,
  \item from one hand-authored feature family to alternative feature dictionaries,
  \item from current witness classes to richer class taxonomies,
  \item from synthetic witnesses to more realistic or mixed-source corpora,
  \item from simple Gaussian evidence models to stronger nonlinear or mixed-mode backends.
\end{enumerate}

What does \emph{not} yet follow automatically is equally important:
\begin{itemize}[leftmargin=1.5em]
  \item current 1D witness-suite success does not prove that every 3D backend is already solved,
  \item current feature-separation plots do not prove that any future class ontology will remain identifiable,
  \item current corpus audits do not remove the need for backend-specific validity and leakage checks after dimensional lift.
\end{itemize}

The practical reading is therefore:
\begin{equation}
\text{current repository status}
\approx
\text{methodology and contract proof}
\;+\;
\text{partial witness evidence},
\end{equation}
not yet
\begin{equation}
\text{final domain-complete classifier stack}.
\end{equation}

That is the correct standard for interpreting the repository today. If future work changes feature sets, class definitions, dimensionality, or sensing backends, the right question is not ``does the old toy still run?'' The right question is:
\begin{quote}
\textit{Do the same corpus, evidence, posterior, and promotion contracts still make the new study legible and auditable?}
\end{quote}

\clearpage
\addcontentsline{toc}{part}{Bibliography}
\begin{thebibliography}{9}

\bibitem{repo_story}
kinematic-classifier-sandbox,
\textit{Repo Story, Workflow Notes, and Companion Methodology Documents},
repository documentation set, 2026.

\bibitem{barshalom}
Yaakov Bar-Shalom, X. Rong Li, and Thiagalingam Kirubarajan,
\textit{Estimation with Applications to Tracking and Navigation},
Wiley, 2001.

\bibitem{bishop}
Christopher M. Bishop,
\textit{Pattern Recognition and Machine Learning},
Springer, 2006.

\bibitem{sutton_barto}
Richard S. Sutton and Andrew G. Barto,
\textit{Reinforcement Learning: An Introduction},
second edition, MIT Press, 2018.

\bibitem{deb}
Kalyanmoy Deb,
\textit{Multi-Objective Optimization Using Evolutionary Algorithms},
Wiley, 2001.

\bibitem{mouret_clune}
Jean-Baptiste Mouret and Jeff Clune,
``Illuminating search spaces by mapping elites,''
\textit{arXiv preprint arXiv:1504.04909}, 2015.

\end{thebibliography}

\end{document}
